\chapter{系统架构与方法}
\label{chap:method}

本章描述智护家系统的总体架构、跨模态注意力融合机制、模态权重设计以及缺失数据处理策略。

\section{系统总体架构}

系统采用``单模态真实分析器 + 学习型融合网络''的架构，如图\ref{fig:architecture}所示。四模态输入分别经特征提取器得到固定维度的特征向量，再经Cross-Modal Attention融合网络进行跨模态交互，最终输出三级风险评估（低/中/高风险）。

\begin{figure}[H]
\centering
\begin{tikzpicture}[
    node distance=0.6cm,
    box/.style={draw, rounded corners, minimum width=2.2cm, minimum height=0.8cm, align=center, font=\small},
    extractor/.style={box, fill=blue!10},
    fusion/.style={box, fill=green!10, minimum width=8cm},
    classifier/.style={box, fill=orange!10, minimum width=6cm},
    input/.style={box, fill=gray!10},
    output/.style={box, fill=red!10, minimum width=4cm},
    arrow/.style={->, >=Stealth, thick}
]

% 输入层
\node[input] (vin) {视频(mp4)};
\node[input, right=of vin] (ain) {音频(wav)};
\node[input, right=of ain] (hin) {生理数据};
\node[input, right=of hin] (min) {用药记录};

% 特征提取层
\node[extractor, below=of vin] (vext) {MediaPipe\\骨架$\to$768维};
\node[extractor, below=of ain] (aext) {librosa\\声学$\to$768维};
\node[extractor, below=of hin] (hext) {统计量\\$\to$256维};
\node[extractor, below=of min] (mext) {规则编码\\$\to$128维};

% 投影层
\node[box, fill=yellow!10, below=1.2cm of vext, minimum width=10cm] (proj) {特征投影层（对齐到512维 + 模态类型嵌入）};

% 融合层
\node[fusion, below=of proj] (fusion) {Cross-Modal Attention $\times$ 4层 + Self-Attention + FFN\\（8 heads, 残差连接+LayerNorm）};

% 分类层
\node[classifier, below=of fusion] (cls) {分类头: 512 $\to$ 256 $\to$ 128 $\to$ 3};

% 输出
\node[output, below=of cls] (out) {风险等级 + 概率分布 + 注意力权重};

% 连接
\draw[arrow] (vin) -- (vext);
\draw[arrow] (ain) -- (aext);
\draw[arrow] (hin) -- (hext);
\draw[arrow] (min) -- (mext);
\draw[arrow] (vext) -- (proj);
\draw[arrow] (aext) -- (proj);
\draw[arrow] (hext) -- (proj);
\draw[arrow] (mext) -- (proj);
\draw[arrow] (proj) -- (fusion);
\draw[arrow] (fusion) -- (cls);
\draw[arrow] (cls) -- (out);

\end{tikzpicture}
\caption{智护家系统总体架构}
\label{fig:architecture}
\end{figure}

各模态特征维度如表\ref{tab:dims}所示。

\begin{table}[H]
\centering
\caption{各模态特征维度}
\label{tab:dims}
\begin{tabular}{lll}
\toprule
模态 & 特征提取方法 & 特征维度 \\
\midrule
视频 & MediaPipe PoseLandmarker 33关键点 & 768 \\
音频 & librosa MFCC + Mel频谱 + 统计量 & 768 \\
生理 & 心率/血氧/血压/步数统计量 & 256 \\
用药 & 应服/已服/漏服/依从率编码 & 128 \\
\bottomrule
\end{tabular}
\end{table}

\section{跨模态注意力机制}
\label{sec:crossattn}

跨模态注意力的核心是让每个模态的特征作为Query，与其他模态的Key/Value进行注意力计算，从而实现模态间信息交互。

\subsection{特征投影与模态嵌入}

给定四个模态特征 $F_v, F_a, F_h, F_m \in \mathbb{R}^{B \times d_i}$（其中 $d_i$ 为各模态原始维度），首先投影到统一维度 $d_{fusion}=512$ 并叠加模态类型嵌入：

\begin{equation}
\tilde{F}_i = W_i F_i + E_{mod}(i), \quad i \in \{v, a, h, m\}
\label{eq:proj}
\end{equation}

其中 $W_i \in \mathbb{R}^{512 \times d_i}$ 为各模态投影矩阵，$E_{mod}(i) \in \mathbb{R}^{512}$ 为可学习的模态类型嵌入。投影层包含LayerNorm、ReLU激活和Dropout（0.2）。

四个模态特征堆叠为序列 $X \in \mathbb{R}^{B \times 4 \times 512}$，送入Cross-Modal Attention层。

\subsection{多头注意力计算}

每层Cross-Modal Attention的计算遵循标准多头注意力\cite{vaswani2017attention}。对于输入序列 $X$，首先通过线性投影得到Query、Key、Value：

\begin{equation}
Q = XW^Q, \quad K = XW^K, \quad V = XW^V
\end{equation}

其中 $W^Q, W^K, W^V \in \mathbb{R}^{512 \times 512}$。注意力计算为：

\begin{equation}
\text{Attention}(Q, K, V) = \text{softmax}\left(\frac{QK^T}{\sqrt{d_k}}\right)V
\label{eq:attention}
\end{equation}

其中 $d_k = 512 / 8 = 64$ 为每个注意力头的维度。8个注意力头并行计算后拼接，经输出投影。

此处Query、Key、Value均来自同一模态序列 $X$，因此实际上是跨模态的自注意力——每个模态token可以attend到所有模态（含自身）。这使得模型能够学习``哪些模态对当前判断更相关''。

\subsection{残差连接与前馈网络}

每层包含残差连接和LayerNorm，遵循Transformer的标准结构：

\begin{equation}
X' = \text{LayerNorm}(X + \text{MultiHead}(X, X, X))
\end{equation}

\begin{equation}
X'' = \text{LayerNorm}(X' + \text{FFN}(X'))
\end{equation}

其中FFN为两层前馈网络（$512 \to 2048 \to 512$，ReLU激活）。该结构堆叠4层。最终取序列的平均池化输出送入分类头。

\subsection{与简单拼接的区别}

MLP拼接（Lite基线）直接将四模态特征concat后过全连接层，模态间无显式交互。Cross-Modal Attention通过注意力权重建模模态间的相关性，是一种更灵活的中间融合。第\ref{chap:exp}章的消融实验（A组）将对比两种融合策略的效果。

\section{模态全局权重}
\label{sec:modweight}

除注意力机制外，网络学习一个可训练的模态权重向量 $\alpha \in \mathbb{R}^4$，经softmax归一化后作为各模态的全局重要性权重：

\begin{equation}
w = \text{softmax}(\alpha) = \frac{e^{\alpha_i}}{\sum_{j=1}^{4} e^{\alpha_j}}, \quad \sum_{i} w_i = 1
\label{eq:modweight}
\end{equation}

融合特征为各模态投影特征的加权求和：

\begin{equation}
F_{global} = \sum_{i} w_i \cdot \tilde{F}_i
\end{equation}

该全局特征与注意力层的平均池化输出拼接后送入分类头。此设计使模型同时具备全局模态加权和细粒度跨模态交互两种能力，且权重 $w$ 可直接用于可解释性展示。

\section{分类头}

分类头为三层全连接网络：

\begin{equation}
\hat{y} = \text{softmax}(\text{MLP}([F_{attn}; F_{global}]))
\end{equation}

其中 $[F_{attn}; F_{global}] \in \mathbb{R}^{1024}$，经 $1024 \to 512 \to 256 \to 3$ 的全连接网络（每层含LayerNorm、ReLU、Dropout 0.3），输出三个风险等级的概率分布。

\section{缺失数据处理}

实际部署中模态数据常缺失（如老人未佩戴手环则无生理数据）。系统采用默认特征填充策略：用训练集中低风险样本的各模态特征均值作为默认值，如表\ref{tab:defaults}所示。

\begin{table}[H]
\centering
\caption{缺失模态的默认特征填充策略}
\label{tab:defaults}
\begin{tabular}{llll}
\toprule
缺失模态 & 默认特征文件 & 含义 & L2范数 \\
\midrule
视频 & normal\_video.npy & 低风险视频特征均值 & 27.27 \\
音频 & silent\_audio.npy & 低风险音频特征均值 & 24.69 \\
生理 & healthy\_baseline.npy & 健康生理基线 & 6.70 \\
用药 & no\_medication.npy & 无用药零向量 & 0.00 \\
\bottomrule
\end{tabular}
\end{table}

此策略保证任意模态缺失时系统仍可推理，且默认值具有真实语义（基于训练集真实特征均值，而非随机噪声）。这使得缺失模态不会引入异常分布，模型能基于可用模态做出合理判断。

\section{模型参数量}

\begin{table}[H]
\centering
\caption{模型参数量对比}
\label{tab:params}
\begin{tabular}{lrr}
\toprule
模型 & 参数量 & 压缩比 \\
\midrule
Full (Cross-Modal Attention) & 8,511,495 & 100\% \\
Lite (MLP Concat) & 3,740,163 & 43.94\% \\
\bottomrule
\end{tabular}
\end{table}

Full模型包含特征投影层（4$\times$）、模态嵌入、4层Cross-Modal Attention（每层含QKV投影+输出投影+FFN）、全局权重和分类头。Lite模型用单层MLP替代Cross-Modal Attention，参数量减少约56\%。
